The most consequential configuration decision in the build was the system clock, and the
reasoning is reproducible from the datasheet, so we set it out rather than leave it as an
implementation footnote.

The ATmega32 derives its UART bit clock by integer division of the system clock. In
standard mode the divisor register is
$\mathrm{UBRR} = \lfloor f_{\mathrm{osc}} / (16 \times \mathrm{BAUD}) + 0.5 \rfloor - 1$,
with the factor 16 replaced by 8 when the double-speed bit U2X is
set~\cite{microchip2020}. Because UBRR is an integer, the achievable baud rate is
quantised, and the quantisation is coarse when $f_{\mathrm{osc}}$ is low. The datasheet's
recommended ceiling on total receiver error for an 8-bit frame with 16$\times$
oversampling is $\pm 2.0\%$~\cite{microchip2020}.

\begin{table}[t]
  \caption{UART divisor error at the factory 1~MHz clock versus the fuse-selected
    8~MHz clock. Rows marked $\dagger$ exceed the datasheet's recommended
    $\pm 2.0\%$ receiver error tolerance~\cite{microchip2020}.}
  \label{tab:baud}
  \small
  \begin{tabular}{@{}llcrrr@{}}
    \toprule
    $f_{\mathrm{osc}}$ & Target & U2X & UBRR & Actual & Error \\
    \midrule
    1~MHz & 38400 & 0 & 1  & 31250 & $-18.6\%^{\dagger}$ \\
    1~MHz & 38400 & 1 & 2  & 41667 & $+8.5\%^{\dagger}$ \\
    \midrule
    8~MHz & 38400 & 0 & 12 & 38462 & $+0.16\%$ \\
    8~MHz & 38400 & 1 & 25 & 38462 & $+0.16\%$ \\
    \bottomrule
  \end{tabular}
\end{table}


Table~\ref{tab:baud} works the arithmetic through. At the factory-default 1~MHz clock,
38400~bps is simply not reachable: the nearest divisor in standard mode yields 31250~bps,
an error of $-18.6\%$, roughly nine times the recommended limit, and double-speed mode
improves this only to $+8.5\%$, still four times over. The 9600~bps case is salvageable but
only in double-speed mode, standard mode giving $-7.0\%$.

Raising the clock to 8~MHz resolves the problem outright: 38400~bps is then reachable with
an error of $+0.16\%$ in either mode, comfortably inside the ceiling. Reliability was
significantly increased by using the internal 8~MHz oscillator and double-speed UART.

Two aspects of how we got there are worth recording, because both are the kind of
constraint that shapes a low-cost build. The change was made by programming the CKSEL fuse
bits rather than by fitting an external crystal, which was unavailable to us at the time;
fuse programming adds no components and requires no board changes, but a mis-set fuse can
leave the part unable to clock itself and recoverable only through high-voltage
programming, so it is not a reversible experiment. And the internal oscillator that makes
this convenient is specified to roughly $\pm 3\%$~\cite{microchip2020}, an order of
magnitude looser than a crystal: adequate for UART framing once the divisor error is this
small, but propagating into every delay-loop-derived interval in the firmware.

The serial path characterised here is wired out but unused by the delivered firmware, which
operates entirely offline per R1. Establishing that 38400~bps is attainable matters because
it keeps the optional uplink of Section~\ref{sec:limitations} available without a further
hardware revision.
